\begin{claim}{D1.3}{Proved}{The untwisted tangent direct images}
For every admissible $c$ in the standard family,
\[
 \boxed{f_*T_X=\OO\oplus\OO(-1),\quad
 R^1f_*T_X=\OO\oplus\OO(-1)^2,\quad R^2f_*T_X=\OO.}   \tag{D.4}
\]
Consequently $(h^0,h^1,h^2,h^3)(X,T_X)=(1,1,1,0)$.
\end{claim}
\begin{proof}
Put $L=\OO_{\mathcal B}(-5)$. The first connecting map of (D.3) is
the holomorphic bundle map
$\kappa:L\to E\otimes G$. In the splittings fixed in D1.2 the
coefficient in $\OO\xi\otimes\OO\gamma$ lies in $H^0(\OO(5))=0$.
The remaining coefficients lie in
$H^0(\OO(4))=\C y$, twice, and $H^0(\OO(3))=\C x$.
Therefore, as a map from $E^*$ to $G$ with its $L$ factor understood,
\[
 \kappa=\begin{pmatrix}0&a y\\ b y&c_0x\end{pmatrix}. \tag{D.5}
\]
We use $c_0$ for this coefficient to distinguish it from the deformation
parameter. The logarithmic period derivative at the cusp has leading
term the invertible matrix $B_0/(2\pi i)$. Hence
$\det\kappa=-ab y^2$ is nonzero there, so $ab\ne0$.
At the order-three point $\tau-\rho$ has order one on the root chart:
$j=1728t$ and the elliptic ramification order of $j$ at $\rho$ is
three. The marked torus period derivative is consequently nonzero
there: the two constant identity columns in the normalized period
matrix exclude an infinitesimal linear-frame change that could
otherwise conceal a nonzero derivative. In (D.5), $y=0$ and
$x\ne0$ at that point, so $c_0\ne0$.
These arguments establish the three nonzero coefficients; they do not
infer them from a global cohomology dimension. Rescaling the two
quotient frames and the source frame makes all three coefficients one.

The next connecting map is cup product with $\kappa$. In suitable
bases it is
\[
 L\otimes G=\OO(-5)\oplus\OO(-6)
 \xrightarrow{\ \delta_1\ }
 E\otimes\det G=\OO(-1)\oplus\OO(-2),\qquad
 \delta_1=\begin{pmatrix}-y&0\\-x&y\end{pmatrix}.       \tag{D.6}
\]
The sign follows by wedging first with $\gamma$ and then with the
second generator of $G$; reversing both sign conventions leaves the
cokernels unchanged. As its determinant is $-y^2$, it is injective as
a morphism of sheaves. The first connecting map is
$\delta_0=(0,y,y,x)^t$ and is a subbundle inclusion, since $x,y$ have
no common zero. Its cokernel is
\[
 \OO\oplus\OO(-1)\oplus\OO(2).                       \tag{D.7}
\]
Indeed one constant linear combination of the two $\OO(-1)$ summands
is unused, and the remaining quotient is given by the exact sequence
$0\to\OO(-5)\xrightarrow{(y,x)}\OO(-1)\oplus\OO(-2)
\xrightarrow{(x,-y)}\OO(2)\to0$.

The cokernel of (D.6) is supported at the order-three root point. There
$x$ is a unit, and elementary row and column operations identify it
with
\[
 Q=\OO(-1)\big|_{2p_1},                               \tag{D.8}
\]
where $2p_1$ means the length-two divisor on the root chart, not the
double coarse point. The generator is the first target component,
and its annihilator is $y^2$. The long direct-image sequence now gives
\[
 R^0\bar f_*\mathcal T=E,\quad
 R^1\bar f_*\mathcal T=\OO\oplus\OO(-1)\oplus\OO(2),
\]
\[
 0\longrightarrow Q\longrightarrow R^2\bar f_*\mathcal T
 \longrightarrow\OO(-6)\longrightarrow0.              \tag{D.9}
\]
This last extension splits, since $\operatorname{Ext}^1(\OO(-6),Q)
=H^1(Q(6))=0$; its splitting is not asserted canonical.

Let $\nu:P=\mathrm dP_6\to W\hookrightarrow X$ be the normalization
immersion. There is an exact sequence
\[
 0\longrightarrow\mathcal T\longrightarrow T_X
 \longrightarrow\nu_*K_P\longrightarrow0.             \tag{D.10}
\]
To verify the quotient, in a chart $W=(z_1\cdots z_r=0)$ the map takes
a vector field to its normal coefficient on each of the $r$ separate
branches. The coefficients lift independently; its kernel is precisely
the vector fields tangent to every branch. Thus the quotient is the
pushforward of the normal line of $\nu$, not $\OO_W(W)$.
The form $d\log x_1\wedge d\log x_2\wedge dq$ becomes a unit times
$dz_1\wedge dz_2\wedge dz_3$ in every unimodular chart and is invariant
under every cusp deck transformation. It trivializes $\nu^*K_X$.
Adjunction for the immersion gives its normal line $K_P$, proving
(D.10). Since $P$ is rational, its $K_P$ cohomology is $(0,0,1)$ by
Serre duality. Finite pushforward preserves that cohomology.

It remains to identify the extension at $p_0$ in degree two. We give
a fibre-sensitive calculation. The kernel of
\[
 \Omega_X^1|_W\longrightarrow\nu_*\Omega_P^1
\]
is $\nu_*\OO_P\otimes N_W^*$. At $xy=0$ its sections are
$a(y,z)y\,dx+b(x,z)x\,dy$; at $xyz=0$ they are
$a(y,z)yz\,dx+b(x,z)xz\,dy+c(x,y)xy\,dz$.
These identifications glue by the common conormal $dq$.
Since $N_W$ and $K_X|_W$ are trivial and $H^0(P,\Omega_P^1)=0$, the
Serre-dual space of $H^2(W,T_X|_W)$ is one-dimensional.
Top-degree base change identifies this group with the residue fibre
of $(R^2\bar f_*T_X)_{p_0}$. Its generic rank is one by (D.9).
A finite module over the discrete valuation ring $\C\{q\}$ of rank
one and residue dimension one is free of rank one. Thus the extension
by $\C_{p_0}$ increases the degree of its line part by $p_0$, without
adding torsion. We obtain the complete root direct images
\[
 R^0\bar f_*T_X=\OO\oplus\OO(-1),\quad
 R^1\bar f_*T_X=\OO\oplus\OO(-1)\oplus\OO(2),
\]
\[
 R^2\bar f_*T_X=\OO(6)\oplus Q.                       \tag{D.11}
\]
On the order-three chart, an invariant coefficient of the generator
of $Q$ must have exponent $2$ modulo $3$, whereas $Q$ only retains
exponents $0,1$. Therefore $\rho_*Q=0$. Applying the exact $\rho_*$
to (D.11) and using
$d(0)=0,d(-1)=d(2)=-1,d(6)=0$ proves (D.4).
Finally Leray over a curve has only columns zero and one. The usual
cohomology of these displayed line bundles gives the four asserted
dimensions. This is a direct-image computation, not an inference from
the tangent Euler characteristic.
\end{proof}
